\section{预处理、编译、汇编和链接各自保留不同信息}

预处理器展开包含关系和条件编译，得到单个翻译单元；编译器解析类型与控制流，
生成目标指令和重定位；汇编器把文本汇编转换为目标文件；链接器合并所有对象，
解决符号并安排最终地址。一个“编译命令”可能驱动全部阶段，但诊断时必须知道
错误属于哪层。

语法错误发生在编译或汇编前端，未定义符号通常在链接时暴露，复位字节位置错误
则可能直到镜像审计或 QEMU 才出现。把每层产物保留下来，可以在最靠近根因的
位置失败。

\section{ELF 目标文件同时承载字节和未完成关系}

可重定位 ELF 中的节还没有最终运行地址。符号表给名字绑定节内偏移，重定位表
记录需要在链接时修补的位置、类型和目标符号。链接器把同类输入节聚合到输出节，
根据链接脚本分配地址，再把符号值代入重定位公式。

\begin{longtable}{L{0.19\textwidth}L{0.26\textwidth}L{0.41\textwidth}}
\toprule
\textbf{信息} & \textbf{典型查看工具} & \textbf{能够确认的事实}\\
\midrule
\endfirsthead
\toprule
\textbf{信息} & \textbf{典型查看工具} & \textbf{能够确认的事实}\\
\midrule
\endhead
ELF 头 & \code{llvm-readelf -h} & 目标机器、类别、端序与文件类型\\
节表 & \code{llvm-readelf -S} & 输入内容分组、大小和属性\\
符号表 & \code{llvm-nm} & 入口、边界和未解析引用\\
重定位 & \code{llvm-readelf -r} & 尚待链接决定的地址字段\\
反汇编 & \code{llvm-objdump -d} & 最终地址上的机器指令\\
裸镜像 & 十六进制审计器 & 平台实际读取的固定偏移字节\\
\bottomrule
\end{longtable}

\section{链接脚本是地址空间的可执行规格}

链接脚本的 MEMORY 区域描述 ROM 容量与属性，SECTIONS 决定输入节的排列，
符号暴露区域边界，ASSERT 把不变量变成链接失败。入口节固定在
\code{0xFFFFF000}，复位节固定在 \code{0xFFFFFFF0}，两者不能仅依赖
源文件顺序偶然到位。

链接器地址计数器表示当前输出地址。对齐表达式可能产生空洞，objcopy 再用
\code{0xFF} 填补。空洞大小同样属于镜像语义：若意外增加一个大节，输出可能
超过 ROM；若缺少固定起点，objcopy 可能生成尺寸完全不同的文件。

\section{从 ELF 转换到 ROM 会主动丢弃信息}

ELF 保存节名、符号、调试信息和地址关系，CPU 复位时不会解析这些元数据。
裸镜像只保留按地址排列的字节。转换是有意的信息压缩，所以调试时同时保留
\code{.elf} 和 \code{.bin}：前者用于解释，后者用于执行。

镜像审计不能只读取 ELF，因为 objcopy 参数、gap fill 和输出路径也可能出错。
项目对最终 \code{firmware.bin} 检查精确尺寸、固定偏移和指令前缀，保证
QEMU 接收的文件与链接意图一致。

\section{C++ 对象模型在裸机上仍然有效}

构造、析构、访问控制、模板和强类型是编译语言机制，不依赖宿主操作系统。
动态分配、异常传播、线程安全静态初始化和程序退出则需要运行时支持。项目可以
使用前一组能力，同时对后一组逐项实现或禁用。

对象生命周期仍须遵守。把一段字节直接解释为 C++ 对象要考虑对齐、有效类型、
平凡复制和构造；磁盘结构不能靠宿主结构体内存布局序列化。系统代码越接近
硬件，越需要类型系统明确哪些转换经过验证。

\section{全局对象初始化需要自研启动协议}

静态存储期对象分为零初始化、常量初始化和动态初始化。链接器把零初始化对象
放入 BSS，入口必须清零；可在编译期完成的常量直接进入镜像；需要执行构造函数
的对象通常由 CRT 遍历初始化数组。

当前基础模块避免依赖动态全局构造。未来确实需要时，内核入口要由链接脚本提供
构造数组边界，并在堆、中断和多线程启动前按顺序调用。析构是否执行也要由内核
生命周期定义，不能默认存在进程退出运行时。

\section{编译器可能生成源码中没有的辅助调用}

大对象复制、除法、栈保护、异常和某些原子操作可能被降低为
\code{memcpy}、\code{__divdi3}、\code{__stack_chk_fail} 或
\code{__cxa_*} 等符号。源码没有显式调用不代表目标文件没有运行时依赖。

项目的 freestanding 符号审计对静态库运行 \code{llvm-nm -u}，任何未知
未解析符号都使构建失败。后续实现 \code{memcpy} 等基础函数时，要把允许列表
收敛到明确自研符号，不能简单忽略全部未定义引用。

\section{内核代码模型影响地址生成方式}

x86-64 指令不能在任意位置直接嵌入完整 64 位地址，编译器根据 code model
选择 RIP 相对、32 位符号扩展或寄存器装载。内核通常位于高半地址，需要与
\code{-mcmodel=kernel}、链接地址和页表布局一致。

\code{-fno-pic} 与 \code{-fno-pie} 让初期内核采用固定链接地址，减少全局
偏移表和动态重定位。未来支持地址随机化时，要设计自举重定位或位置无关入口，
不能只切换编译选项。

\section{头源分离同时控制依赖与 ABI}

头文件公开稳定类型与函数，源文件隐藏实现。改动私有算法不应改变上层编译
依赖，改动公开类型布局则是 ABI 变化。裸机模块没有动态库版本协商，因此同一
镜像中的组件必须由同一提交共同构建。

硬件寄存器细节放在驱动私有实现，公共接口暴露领域动作和状态；页表位掩码放在
内存模块，调度器只看到地址空间对象。这样的分层防止魔法数字从底层泄漏到
系统策略。

\section{ROM 入口逐步建立 C++ 可运行条件}

当前 v0.2 已建立 16 位栈、UART 和经过校验的磁盘加载，但还没有调用 C++。
通往 C++ 内核的完整依赖链是：

\begin{enumerate}[label=\arabic*.]
  \item ROM 入口关闭未配置中断并设置确定方向标志。
  \item 磁盘驱动把后续阶段读入经过验证的 RAM；v0.2 已完成。
  \item GDT 与页表使处理器进入 64 位长模式。
  \item ELF 装载器复制段、清零 BSS 并检查入口。
  \item 汇编适配层选择对齐栈并传入 BootInfo。
  \item C++ 入口建立必要运行时，再初始化内核子系统。
\end{enumerate}

每一步只消费上一阶段已经证明的状态。若直接在 16 位复位入口调用 64 位 C++，
编译器假设的寄存器、指令模式、栈对齐和地址模型全部不成立。

\section{可重复构建把产物绑定到源状态}

同一提交与工具版本应产生同一 ROM 哈希。构建脚本不能把当前时间、随机路径或
未初始化宿主内存带入镜像。产物哈希使测试日志、发布记录和调试现场指向唯一
字节集合。

可重复不等于跨任意编译器版本字节相同。工具升级可能改变合法指令选择，因此
发布记录还要保存工具版本；升级后重新审计尺寸、反汇编和运行行为。
